
\SetTituloActual{Macetas biodegradables de hojas recicladas}

\Inicio{ARTÍCULO DE INVESTIGACIÓN}{Uso de hojas de árbol para la creación de macetas biodegradables}{Zamora Perez Lesly María, Cota Velarde Xochitl Noemí, Gonzalez Sandoval Blanca Nelly}

\Resumen{Este proyecto analiza la problemática de la acumulación de hojas en zonas urbanas, como Mazatlán, donde estos residuos orgánicos pueden provocar obstrucción de drenajes, inundaciones, malos olores y afectaciones a la salud. Además, se considera el impacto ambiental del uso excesivo de macetas de plástico, las cuales generan contaminación por su lenta degradación.\\
Como alternativa, se propone la elaboración de macetas biodegradables a partir de hojas secas reutilizadas, aplicando el enfoque de economía circular. La metodología consistió en triturar las hojas, mezclarlas con agua y harina como aglutinante, moldearlas y secarlas al sol para obtener un producto funcional. Se llevó a cabo un experimento comparativo con macetas de plástico para evaluar su efectividad en el crecimiento vegetal.\\
Diversos antecedentes respaldan la viabilidad de esta propuesta, demostrando que los materiales vegetales pueden sustituir al plástico sin afectar el crecimiento de las plantas. Los análisis indican que las macetas biodegradables favorecen el desarrollo radicular y aumentan la supervivencia tras el trasplante. Aunque los resultados mostraron una tendencia favorable, no fueron estadísticamente significativos (p > 0.05), lo que sugiere la necesidad de estudios con muestras más grandes.\\
En conclusión, este proyecto ofrece una solución sustentable que reduce residuos orgánicos y plásticos, contribuye al cuidado del medio ambiente y promueve prácticas ecológicas en la comunidad.}{Hojas secas, macetas biodegradables, sustentabilidad, residuos orgánicos, reducción de plásticos.}

\begin{multicols}{2}

\section{Introducción}

En zonas urbanas con abundante vegetación, como Mazatlán, es frecuente la acumulación de hojas de árboles, especialmente durante ciertas temporadas del año. Aunque se trata de residuos naturales, su acumulación puede generar problemas como la obstrucción de alcantarillas, inundaciones, malos olores y afectaciones a la salud pública \citep{crowther2015mapping}. Esta problemática se ve agravada por la falta de estrategias de gestión adecuadas para este tipo de residuos orgánicos \citep{fao2019food}.

Paralelamente, el uso excesivo de plásticos, particularmente en productos de uso cotidiano como las macetas, representa un grave problema ambiental debido a su lenta degradación y su impacto en los ecosistemas \citep{unep2018singleuse}. La producción y disposición de plásticos contribuyen significativamente a la contaminación del suelo y del agua, afectando la calidad de vida de las comunidades y los ecosistemas locales.

Ante esta doble problemática, surge la necesidad de replantear el manejo de los residuos, considerándolos como un recurso potencial bajo un enfoque de economía circular. En este sentido, el presente estudio propone la utilización de hojas secas para la elaboración de macetas biodegradables, como una alternativa ecológica que permita reducir tanto la acumulación de residuos orgánicos como el uso de plásticos \citep{fuentes2021development}. Esta propuesta se enmarca dentro de las estrategias de desarrollo sostenible, promoviendo el aprovechamiento de recursos naturales y la implementación de soluciones sostenibles en el entorno urbano.

\section{Materiales y métodos}

\subsection{Materiales y procedimiento}

Para la elaboración de las macetas biodegradables se utilizaron hojas secas recolectadas en áreas urbanas de Mazatlán. El proceso comenzó con el triturado manual de las hojas hasta obtener una textura fina y homogénea. Posteriormente, se añadió agua y 200 g de harina de trigo como aglutinante natural, junto con una pequeña cantidad de vinagre para ayudar en la conservación del material. La mezcla se homogeneizó hasta formar una pasta consistente y maleable.

Esta pasta se vertió en moldes con la forma deseada para macetas, compactándola adecuadamente para asegurar su estructura. Finalmente, las macetas fueron expuestas a la luz solar directa durante varios días hasta su completo secado y endurecimiento, lo que les proporcionó la rigidez necesaria para su uso como contenedores de plantas.

\subsection{Diseño experimental y análisis estadístico}

Para evaluar la efectividad de las macetas biodegradables, se diseñó un experimento comparativo con dos grupos experimentales: 6 macetas de plástico (grupo control) y 6 macetas biodegradables (grupo experimental). Se sembró la misma especie vegetal en ambos grupos, manteniendo condiciones ambientales y de riego uniformes.

Las variables medidas fueron: número de plantas emergidas (tasa de germinación), altura de la planta (medida con una regla desde la base del tallo hasta el ápice) y porcentaje de raíz seca (calculado después de extraer las plantas y secar las raíces a temperatura ambiente). La medición de la altura se realizó a los 15 días después de la siembra, mientras que la evaluación de las raíces se efectuó al final del experimento, a los 30 días.

Los datos obtenidos fueron procesados mediante una prueba t de Student para muestras independientes, con un nivel de significancia de p < 0.05, utilizando el software estadístico JASP. Este análisis permitió comparar las medias de ambos grupos y determinar si las diferencias observadas eran estadísticamente significativas.

\section{Resultados}

Los resultados del experimento mostraron un mejor desempeño inicial de las plantas cultivadas en macetas biodegradables en comparación con las de plástico. La \autoref{tab:resultados_dos_grupos} presenta los valores obtenidos para cada variable medida.

\begin{table}[H]

\caption{Resultados de los dos grupos}
\label{tab:resultados_dos_grupos}

\begin{threeparttable}

{\centering\small\setstretch{1.35}\rowcolors{2}{Grey}{White}
\begin{tabular}{M{0.3\linewidth}M{0.275\linewidth}M{0.275\linewidth}}\hline

    \theader{Variable} & \theader{Plástico (MP)} & \theader{Biodegradable (MB)} \\ \hline
    
    Número de plantas emergidas & 2 de 6       & 3 de 6 \\
    Altura de la planta (cm)    & 1 y 2        & 2.5, 2.8 y 3 \\
    Altura promedio (cm)        & 1.5          & 2.77 \\
    Raíz seca                   & 20\% y 40\%  & 10\%, 20\% y 20\% \\
    Promedio de raíz seca       & 30\%         & 16.7\% \\
    Raíces deformadas           & 0\%          & 0\% \\
    
    \hline
\end{tabular}}

\begin{tablenotes}\tiny % ex: \tnote{1} ... \item[1]
    \item[] 
\end{tablenotes}

\end{threeparttable}
\end{table}

Como se observa, la altura promedio de las plantas fue mayor en las macetas biodegradables (2.77 cm) que en las de plástico (1.5 cm). Asimismo, el porcentaje de raíz seca fue menor en el grupo experimental (16.7\%) en comparación con el control (30\%). La tasa de germinación también fue superior en las macetas biodegradables (50\%) frente a las de plástico (33\%). No se observaron deformaciones radiculares en ninguno de los tratamientos.

El análisis estadístico mediante la prueba t de Student indicó que las diferencias encontradas no fueron significativas (altura: p ≈ 0.08; raíz seca: p ≈ 0.12). La falta de significancia estadística probablemente se debe al tamaño reducido de la muestra (n=6 por grupo), lo que limita el poder estadístico del análisis.

\section{Discusión}

Los resultados obtenidos sugieren que las macetas biodegradables elaboradas a partir de hojas secas pueden proporcionar condiciones favorables para el crecimiento vegetal, particularmente en el desarrollo radicular. Aunque no se alcanzó significancia estadística, las tendencias observadas indican ventajas en términos de retención de humedad y aireación, factores críticos para el desarrollo inicial de las plantas.

Estos hallazgos coinciden con estudios previos que destacan el potencial de los materiales vegetales como sustitutos del plástico en aplicaciones agrícolas. Por ejemplo, \citet{fuentes2021development} encontraron que las macetas elaboradas con residuos agroindustriales presentaban una mayor tasa de germinación y un mejor desarrollo radicular en comparación con las macetas de plástico. De manera similar, \citet{jiffygroup} ha comercializado exitosamente macetas biodegradables a base de fibras de celulosa, demostrando su viabilidad a escala industrial.

La reducción en el porcentaje de raíz seca observada en las macetas biodegradables (16.7\% vs 30\%) sugiere una mejor salud radicular, lo que podría traducirse en una mayor supervivencia de las plantas tras el trasplante. Este resultado es consistente con lo reportado por \citet{roussny2023macetas}, quienes encontraron que las raíces en macetas biodegradables presentaban menor estrés y mejor desarrollo.

Además, la reutilización de hojas contribuye a la reducción de problemas ambientales urbanos, como la obstrucción de drenajes y la acumulación de residuos, tal como se ha documentado en iniciativas similares como Releaf Paper \citep{ekosnegocios2025releaf}. Esta doble función -ambiental y agronómica- refuerza la pertinencia de esta propuesta dentro de un enfoque de economía circular.

No obstante, es importante reconocer las limitaciones del presente estudio. El tamaño reducido de la muestra (n=6 por grupo) limita la capacidad para detectar diferencias significativas y la generalización de los resultados. Futuros estudios deberían considerar un mayor número de réplicas para aumentar el poder estadístico y confirmar las tendencias observadas.

\section{Conclusión}

Las macetas biodegradables elaboradas a partir de hojas secas representan una alternativa viable frente a las macetas de plástico. A pesar de que los resultados no fueron estadísticamente significativos, se observó una tendencia favorable en el crecimiento vegetal y la calidad radicular, lo que sugiere que este material podría ser una opción sustentable para la producción de contenedores de plantas.

Esta propuesta contribuye a la reducción de residuos orgánicos y plásticos, promoviendo prácticas sostenibles y el aprovechamiento de recursos naturales dentro de un enfoque de economía circular. Sin embargo, se recomienda realizar estudios con una mayor cantidad de macetas (n > 30 por grupo) para reforzar los resultados del análisis estadístico y poder establecer conclusiones más robustas sobre la efectividad de estas macetas en comparación con las de plástico.

% ------------------------------------------------------ %

    \bibliographystyle{apalike-ejor}
    \bibliography{art_01/bib}

% ------------------------------------------------------ %

\end{multicols}